\chapter{Custom-Gradient Wrappers for HMC}
\label{ch:bf-custom-gradient-wrappers}

A custom gradient is safe for HMC only if it returns the derivative of the same
scalar log target used in the Metropolis correction.  A fast gradient of a
nearby surrogate changes the dynamics and can bias the chain unless the sampler
is explicitly corrected.

\section{Same-Scalar Contract}
\label{sec:bf-custom-gradient-same-scalar}

Let the production value path return the full sampler-coordinate posterior target
\begin{equation}
\label{eq:bf-custom-primal-target}
  \mathcal{L}(u)
  =
  \ell(T(u);y_{1:T})
  + \log p(T(u))
  + \log\left|\det DT(u)\right|,
\end{equation}
where the likelihood term uses the same structural provider, masks, factor
branch, regularization policy, and parameter transforms declared in
Chapters~\ref{ch:bf-structural-derivatives}--\ref{ch:bf-factor-derivatives}.
When a chapter later writes a likelihood-only derivative such as
\(\nabla_\theta \ell\), that object should be read as one component of the full
sampler-coordinate target gradient, not as the entire HMC target by itself.  The
custom-gradient path is valid only if it returns
\begin{equation}
\label{eq:bf-custom-gradient-target}
  g_i(u)=\frac{\partial \mathcal{L}(u)}{\partial u_i}.
\end{equation}
If the value path evaluates a floored spectral covariance $P_+$, the gradient
must be the derivative of that implemented value path.  If the value path uses a
jittered innovation covariance $S_t+\delta I$, the gradient must include the
same jitter policy or clearly report that the parameter point is outside the
certified smooth branch.  The wrapper must not evaluate one scalar and return the
gradient of another.

BayesFilter's custom-gradient wrapper policy is:
\begin{enumerate}[label=(\roman*)]
  \item compute the primal log likelihood with the production backend;
  \item return an analytic score that has been validated against the same
    primal value on controlled cases;
  \item expose diagnostics for finite failures and backend fallbacks;
  \item keep tensor shapes static inside compiled sampler paths;
  \item use Hessians for mass matrices or diagnostics, not inside every HMC
    leapfrog step unless explicitly justified.
\end{enumerate}

This policy reflects the MacroFinance motivation: nested tape Hessians through
full Kalman scans are useful as small-model validation oracles, but they are not
the intended large-model production path.

\section{Gradient Callback Contract}
\label{sec:bf-custom-gradient-callback}

A custom-gradient wrapper should expose a contract equivalent to the following
mathematical tuple:
\[
  \operatorname{value\_and\_score}(u)
  =
  \left(
    \mathcal{L}(u),
    \nabla_u\mathcal{L}(u),
    d(u)
  \right),
\]
where $d(u)$ is a diagnostic record.  The diagnostic record must include:
\begin{itemize}
  \item the structural provider and parameter transform identifiers;
  \item the active observation-mask policy;
  \item the covariance/factor backend used by the value path;
  \item any jitter, nugget, PSD projection, rank truncation, active spectral
    floor, pivoting, sign correction, or fallback branch;
  \item the minimum relevant Cholesky pivot, QR diagonal, or spectral gap;
  \item finite-value and finite-gradient status.
\end{itemize}
The gradient callback should return a failure status rather than silently
replacing the derivative with zeros, prior gradients, or a different backend.
Emergency fallbacks can be useful for optimization diagnostics, but HMC must
know whether it is using the certified derivative of the returned scalar target.
If the wrapper deliberately returns a finite invalid-region fallback scalar, then
the only HMC-safe fallback gradient is the exact derivative of that fallback
scalar.  In particular, a constant low-density fallback must have zero gradient;
a nonzero surrogate direction defines a different target unless accompanied by an
explicit correction scheme.

\section{Hessian and Observed-Information Path}
\label{sec:bf-custom-gradient-hessian-path}

The HMC hot path needs values and gradients.  Hessians are usually used outside
the leapfrog loop for mass-matrix initialization, local identification, and
curvature diagnostics.  If a wrapper exposes a Hessian, it should target
\[
  H_{ij}(\psi)
  =
  \frac{\partial^2\mathcal{L}(\psi)}
       {\partial\psi_i\,\partial\psi_j}
\]
for the same scalar in \eqref{eq:bf-custom-primal-target}.  The observed
information is $-H(\psi)$, possibly plus a sign-adjusted convention for priors
if the likelihood-only and posterior targets are both reported.  A Hessian path
that omits prior curvature, transformed-parameter curvature, or jitter/floor
branch terms must label itself as a partial diagnostic, not as posterior
curvature.

The recommended production pattern is:
\begin{enumerate}[label=(\roman*)]
  \item value-plus-score path for MAP and HMC;
  \item analytic or custom Hessian path for mode diagnostics and mass-matrix
    construction;
  \item small-model autodiff Hessian path as an oracle, not as the large-model
    default.
\end{enumerate}

\section{Static-Shape and Compilation Contract}
\label{sec:bf-custom-gradient-static-shapes}

Compiled HMC paths require stable shapes.  Parameter dimension, state dimension,
observation dimension, maximum active observation size, and the layout of
diagnostic records should be static across leapfrog calls.  Missing data should
be represented by fixed-shape masks or predeclared selection patterns rather
than by dynamically changing tensor ranks.  Fixed sigma-point rules must keep the
number and ordering of points stable.  Factor backends must report branch events
without changing the differentiable return signature.

This static-shape rule is an implementation constraint, but it has mathematical
content: changing the set of coordinates, points, or factors during a compiled
scan can turn a smooth derivative contract into a piecewise branch derivative
without making the branch visible to the sampler.

\section{HMC Safety Labels}
\label{sec:bf-custom-gradient-hmc-labels}

BayesFilter should distinguish the following labels:
\begin{itemize}
  \item \textbf{value-correct}: the scalar likelihood or posterior value matches
    a reference value path;
  \item \textbf{gradient-correct}: the returned score matches finite
    differences or autodiff of that same scalar on certified cases;
  \item \textbf{curvature-correct}: the Hessian or observed information matches
    the same scalar and sign convention;
  \item \textbf{sampler-usable}: short HMC runs have finite log probabilities,
    finite gradients, and acceptable diagnostic telemetry;
  \item \textbf{production-gated}: compiled/eager parity, branch diagnostics,
    mask tests, stress tests, and representative workload tests all pass.
\end{itemize}
A finite HMC run is not by itself a gradient proof.  Conversely, a correct
gradient on small smooth systems does not certify spectral floors, rank changes,
or DSGE deterministic-completion branches unless those regimes have been
validated explicitly.
